\chapter{Complexity — Zero P Eval}
%%% generated by the blueprint pipeline from TCSlib.Complexity.Randomized.ZeroPEval
\section{Overview}

Declarations of the Lean module \texttt{TCSlib.Complexity.Randomized.ZeroPEval}.

\section{Declarations}

\begin{definition}[Gate value modulo an integer]
\lean{Randomized.PIT.gateVal}
\statementsource{Vad12}{pdf-bf2ce5422088-p016-b001}
\label{Randomized.PIT.gateVal}
\leanok
\uses{Randomized.AlgGate}
Fix a natural number $n$, and let $k$ be a natural number, $a$ a function from natural
numbers to natural numbers, and $vs$ a list of natural numbers. The value of an input
gate $x_i$ is $a(i)$ modulo $k$, and the value of the constant-one gate is $1$ modulo
$k$. An addition, subtraction, or multiplication gate with indices $j,l$ has value,
respectively, the sum, the sum of the left value and $k$ minus the right value, or the
product of the entries of $vs$ at those indices, all modulo $k$; an out-of-range list
lookup is interpreted as zero.
\end{definition}

\begin{definition}[Iterated gate values]
\lean{Randomized.PIT.evalVals}
\statementsource{Vad12}{pdf-bf2ce5422088-p016-b001}
\label{Randomized.PIT.evalVals}
\leanok
\uses{Randomized.AlgGate, Randomized.PIT.gateVal}
Let $n,k$ be natural numbers, let $a:\mathbb N\to\mathbb N$, let $gs$ be a list of
algebraic gates with $n$ inputs, and let $vs$ be a list of natural numbers. Define the
evaluated list by processing the gates in order, appending to the current list the value
of each gate computed modulo $k$ from $a$ and the current list, starting with $vs$.
\end{definition}

\begin{theorem}[Defaulted list lookup and casting]
\lean{Randomized.PIT.getD_map_natCast}
\label{Randomized.PIT.getD_map_natCast}
\leanok
\difficulty{4}
For any modulus $k\in\mathbb{N}$, list $v$ of natural numbers, and index
$j\in\mathbb{N}$, looking up the $j$th entry of the list obtained by casting every entry
of $v$ into $\mathbb{Z}/k\mathbb{Z}$, using $0$ if the index is out of range, gives the
cast into $\mathbb{Z}/k\mathbb{Z}$ of the $j$th entry of $v$, again using $0$ if the
index is out of range.
\end{theorem}

\begin{theorem}[Bound on a Defaulted List Entry]
\lean{Randomized.PIT.getD_le}
\label{Randomized.PIT.getD_le}
\leanok
\difficulty{4}
Let $k$ be a natural number and let $vs$ be a list of natural numbers such that every
entry of $vs$ is strictly less than $k$. For every natural number $j$, the entry at
position $j$, using $0$ when the position is outside the list, is at most $k$.
\end{theorem}

\begin{theorem}[Gate Values Are Bounded by the Modulus]
\lean{Randomized.PIT.gateVal_lt}
\label{Randomized.PIT.gateVal_lt}
\leanok
\uses{Randomized.AlgGate, Randomized.PIT.gateVal}
\difficulty{4}
Let $n$ be a natural number, let $k$ be a positive natural number, let
$a:\mathbb{N}\to\mathbb{N}$, let $vs$ be a list of natural numbers, and let $g$ be an
arithmetic gate on $n$ inputs. The value of $g$ computed from $a$ and $vs$ modulo $k$ is
strictly less than $k$.
\end{theorem}

\begin{theorem}[Gate evaluation modulo an integer]
\lean{Randomized.PIT.natCast_gateVal}
\label{Randomized.PIT.natCast_gateVal}
\leanok
\uses{Randomized.AlgGate, Randomized.AlgGate.eval, Randomized.PIT.gateVal, Randomized.PIT.getD_le, Randomized.PIT.getD_map_natCast}
\difficulty{4}
Let $n,k\in\mathbb{N}$, let $a:\mathbb{N}\to\mathbb{N}$, and let $v$ be a list of
natural numbers such that every entry of $v$ is less than $k$. For every arithmetic gate
$g$ with $n$ possible input variables, evaluating its natural-number gate value modulo
$k$ gives the same element of $\mathbb{Z}/k\mathbb{Z}$ as evaluating the gate over
$\mathbb{Z}/k\mathbb{Z}$ on the input values obtained by reducing $a(i)$ modulo $k$.
\end{theorem}

\begin{theorem}[Evaluation modulo a positive integer]
\lean{Randomized.PIT.map_evalVals}
\label{Randomized.PIT.map_evalVals}
\leanok
\uses{Randomized.AlgGate, Randomized.AlgGate.eval, Randomized.PIT.evalVals, Randomized.PIT.gateVal, Randomized.PIT.gateVal_lt, Randomized.PIT.natCast_gateVal}
\difficulty{5}
Let $k$ be a positive natural number, let $a$ assign a natural number to each input
index, and let $gs$ be a finite list of algebraic gates. For every initial list $vs$ of
natural numbers whose entries are all less than $k$, mapping the list produced by
evaluating $gs$ modulo $k$ from $vs$ into $\mathbb Z/k\mathbb Z$ gives the same list as
evaluating the gates in that ring from the mapped initial list; moreover, every entry of
the natural-number evaluation list is less than $k$.
\end{theorem}

\begin{theorem}[Evaluation agrees with ring evaluation]
\lean{Randomized.PIT.natCast_evalVals}
\label{Randomized.PIT.natCast_evalVals}
\leanok
\uses{Randomized.AlgCircuit.run, Randomized.AlgGate, Randomized.PIT.evalVals, Randomized.PIT.map_evalVals}
\difficulty{1}
Let $k$ be a positive natural number, let $a$ assign a natural number to each input
index, and let $gs$ be a finite list of algebraic gates. Evaluating the gates modulo $k$
from the empty list and mapping the resulting values into $\mathbb Z/k\mathbb Z$ gives
the same list as running those gates in $\mathbb Z/k\mathbb Z$ with each input assigned
the residue of its value under $a$.
\end{theorem}

\begin{theorem}[Range of evaluated gate values]
\lean{Randomized.PIT.lt_of_mem_evalVals}
\label{Randomized.PIT.lt_of_mem_evalVals}
\leanok
\uses{Randomized.AlgGate, Randomized.PIT.evalVals, Randomized.PIT.map_evalVals}
\difficulty{1}
Let $k$ be a positive natural number, let $a$ assign a natural number to each input
index, and let $gs$ be a finite list of algebraic gates. Every value in the list
obtained by evaluating these gates modulo $k$, starting from the empty list of prior
values, is strictly less than $k$.
\end{theorem}

\begin{theorem}[Length of the evaluated value list]
\lean{Randomized.PIT.length_evalVals}
\label{Randomized.PIT.length_evalVals}
\leanok
\uses{Randomized.AlgGate, Randomized.PIT.evalVals, Randomized.PIT.gateVal}
\difficulty{4}
For every modulus $k\in\mathbb{N}$, assignment $a:\mathbb{N}\to\mathbb{N}$, finite list
$gs$ of arithmetic gates, and initial list $vs$ of natural numbers, evaluating the gates
in order and appending their values increases the list length by exactly the number of
gates: the resulting list has length $|vs|+|gs|$.
\end{theorem}

\begin{definition}[One circuit-evaluation step]
\lean{Randomized.PIT.absStep}
\label{Randomized.PIT.absStep}
\leanok
\uses{Randomized.AlgGate.parse, Randomized.AlgGate.refsBelow, Randomized.PIT.gateVal}
For a number of input variables $n$, a modulus $k$, and an assignment
$a:\mathbb{N}\to\mathbb{N}$, define a transition on a pair consisting of a bit list and
a list of gate values as follows. If the bit list begins with true, parse the remaining
bits as a gate and a leftover bit list; when parsing succeeds and every referenced node
index is smaller than the current number of gate values, return the leftover bits
together with the old values followed by the gate’s value modulo $k$. If parsing fails
or a reference is not below that number, return the empty bit list and retain the
values; an empty input stays empty, the singleton false list stays unchanged, and any
other list beginning with false returns the empty bit list, always retaining the values.
\end{definition}

\begin{theorem}[Reconstruction from a successful parse]
\lean{Randomized.PIT.parse_eq_some}
\label{Randomized.PIT.parse_eq_some}
\leanok
\uses{BoolCircuit.decodeNat, BoolCircuit.encodeNat, Complexity.BinArith.decodeNat_eq, Complexity.BinArith.eq_encodeNat_of_hasFalse, Randomized.AlgGate, Randomized.AlgGate.encode, Randomized.AlgGate.parse, Randomized.AlgGate.parseBin}
\difficulty{6}
Let $n\in\mathbb{N}$, let $l,l'$ be bit strings, and let $g$ be an arithmetic gate with
$n$ possible input variables. If parsing $l$ with the gate parser for $n$ variables
returns the gate $g$ and remainder $l'$, then $l$ is the encoding of $g$ followed by
$l'$.
\end{theorem}

\begin{theorem}[Iteration of gate decoding]
\lean{Randomized.PIT.iterate_absStep_encodeGates}
\label{Randomized.PIT.iterate_absStep_encodeGates}
\leanok
\uses{Randomized.AlgCircuit.encodeGates, Randomized.AlgGate, Randomized.AlgGate.parse_encode, Randomized.AlgGate.refsBelow, Randomized.PIT.absStep, Randomized.PIT.evalVals, Randomized.PIT.gateVal}
\difficulty{5}
Fix an input width $n$, a modulus $k\in\mathbb{N}$, an assignment
$a:\mathbb{N}\to\mathbb{N}$, a list $g_1,\ldots,g_m$ of arithmetic gates on $n$ inputs,
and an initial list $v$ of natural numbers. Suppose that each gate's references point
only to earlier values: for every $i<m$, the gate at position $i$ has all references
below $|v|+i$. Then applying the decoding step $m$ times to the encoding of the gate
list and the initial values yields the terminal marker together with the list obtained
by evaluating the gates in order, appending each gate's value to the values already
computed.
\end{theorem}

\begin{theorem}[Parsing an Iterated Absolute Step]
\lean{Randomized.PIT.exists_of_iterate_absStep}
\label{Randomized.PIT.exists_of_iterate_absStep}
\leanok
\uses{Randomized.AlgCircuit.encodeGates, Randomized.AlgGate, Randomized.AlgGate.parse, Randomized.AlgGate.refsBelow, Randomized.PIT.absStep, Randomized.PIT.evalVals, Randomized.PIT.gateVal, Randomized.PIT.parse_eq_some}
\difficulty{6}
Fix a natural number $n$, a modulus $k$, and a function $a:\mathbb{N}\to\mathbb{N}$. For
every natural number $T$, bit string $w$, and lists $vs,vs'$ of natural numbers, if
iterating the absolute parsing step $T$ times from $(w,vs)$ yields $( [0],vs')$, then
there is a list $gs$ of arithmetic gates on $n$ inputs such that $w$ is its gate
encoding, every gate at position $i$ refers only to earlier positions in the
concatenation of $vs$ and the gates, and $vs'$ is the result of evaluating those gates
modulo $k$ starting with $vs$.
\end{theorem}

\begin{definition}[Bit at a position]
\lean{Randomized.PIT.tagBit}
\label{Randomized.PIT.tagBit}
\leanok
For a bit list $w$ and an index $i\in\mathbb{N}$, let the tagged bit be the bit at
position $i$ in $w$, counting from zero; if that position is beyond the end of the list,
its value is false.
\end{definition}

\begin{definition}[Bit-list body]
\lean{Randomized.PIT.bodyW}
\label{Randomized.PIT.bodyW}
\leanok
For a bit list $w$, define its body to be the list obtained by removing its first four
bits, or all its bits if it has fewer than four.
\end{definition}

\begin{definition}[First argument segment]
\lean{Randomized.PIT.arg1W}
\label{Randomized.PIT.arg1W}
\leanok
\uses{Complexity.BinArith.uHead, Randomized.PIT.bodyW}
For a bit list $w$, define its first argument segment to be the result of applying the
operation $uHead$ to the body of $w$.
\end{definition}

\begin{definition}[Middle bit-list segment]
\lean{Randomized.PIT.midW}
\label{Randomized.PIT.midW}
\leanok
\uses{Complexity.BinArith.uRest, Randomized.PIT.bodyW}
For a bit list $w$, define its middle segment to be the result of applying the operation
$uRest$ to the body of $w$.
\end{definition}

\begin{definition}[Second argument segment]
\lean{Randomized.PIT.arg2W}
\label{Randomized.PIT.arg2W}
\leanok
\uses{Complexity.BinArith.uHead, Randomized.PIT.midW}
For a bit list $w$, define its second argument segment to be the result of applying the
operation $uHead$ to its middle segment.
\end{definition}

\begin{definition}[Remaining bit-list segment]
\lean{Randomized.PIT.rest2W}
\label{Randomized.PIT.rest2W}
\leanok
\uses{Complexity.BinArith.uRest, Randomized.PIT.midW}
For a bit list $w$, define its remaining segment to be the result of applying the
operation $uRest$ to its middle segment.
\end{definition}

\begin{definition}[Variable-gate tag]
\lean{Randomized.PIT.isVarW}
\label{Randomized.PIT.isVarW}
\leanok
\uses{Randomized.PIT.tagBit}
For a bit list $w$, define its variable-gate tag to be true exactly when the bits at
positions $1$, $2$, and $3$ are all false, interpreting any position beyond the end of
the list as false.
\end{definition}

\begin{definition}[One-operation tag test]
\lean{Randomized.PIT.isOneW}
\label{Randomized.PIT.isOneW}
\leanok
\uses{Randomized.PIT.tagBit}
For a Boolean list $w$, define the one-operation test to return true exactly when the
tests at positions $1$ and $2$ are false and the test at position $3$ is true.
\end{definition}

\begin{definition}[Addition-operation tag test]
\lean{Randomized.PIT.isAddW}
\label{Randomized.PIT.isAddW}
\leanok
\uses{Randomized.PIT.tagBit}
For a Boolean list $w$, define the addition-operation test to return true exactly when
the tests at positions $1$ and $3$ are false and the test at position $2$ is true.
\end{definition}

\begin{definition}[Subtraction-operation tag test]
\lean{Randomized.PIT.isSubW}
\label{Randomized.PIT.isSubW}
\leanok
\uses{Randomized.PIT.tagBit}
For a Boolean list $w$, define the subtraction-operation test to return true exactly
when the test at position $1$ is false and the tests at positions $2$ and $3$ are true.
\end{definition}

\begin{definition}[Multiplication-operation tag test]
\lean{Randomized.PIT.isMulW}
\label{Randomized.PIT.isMulW}
\leanok
\uses{Randomized.PIT.tagBit}
For a Boolean list $w$, define the multiplication-operation test to return true exactly
when the test at position $1$ is true and the tests at positions $2$ and $3$ are false.
\end{definition}

\begin{definition}[Completion test]
\lean{Randomized.PIT.isDone}
\label{Randomized.PIT.isDone}
\leanok
\uses{Complexity.isNilB}
For a Boolean list $w$, define the completion test to return true if either $w$ is nil
according to the nil test, or its tail after removing the first entry is nil according
to that test and its first entry, defaulting to false when absent, is false.
\end{definition}

\begin{definition}[Validity test for a word]
\lean{Randomized.PIT.okW}
\label{Randomized.PIT.okW}
\leanok
\uses{Complexity.BinArith.hasFalse, Randomized.PIT.arg1W, Randomized.PIT.arg2W, Randomized.PIT.bodyW, Randomized.PIT.isAddW, Randomized.PIT.isMulW, Randomized.PIT.isOneW, Randomized.PIT.isSubW, Randomized.PIT.isVarW, Randomized.PIT.midW}
For natural numbers $n$ and $i$ and a Boolean list $w$, define the validity test to
require that the first entry of $w$, defaulting to false when absent, is true. It also
requires either that $w$ passes the variable test, its body contains a false entry, and
the length of its first argument is less than $n$; that it passes the one-operation
test; or that it passes at least one of the addition, subtraction, or multiplication
tests, its body and middle part each contain a false entry, and the lengths of both
arguments are less than $i$.
\end{definition}

\begin{definition}[Residual part of a word]
\lean{Randomized.PIT.restW}
\label{Randomized.PIT.restW}
\leanok
\uses{Randomized.PIT.bodyW, Randomized.PIT.isOneW, Randomized.PIT.isVarW, Randomized.PIT.midW, Randomized.PIT.rest2W}
For a Boolean list $w$, define its residual part to be its middle part if it passes the
variable test; otherwise, if it passes the one-operation test, take its body; otherwise
take its second residual part.
\end{definition}

\begin{definition}[Modular word evaluation]
\lean{Randomized.PIT.valN}
\label{Randomized.PIT.valN}
\leanok
\uses{Randomized.PIT.arg1W, Randomized.PIT.arg2W, Randomized.PIT.isAddW, Randomized.PIT.isOneW, Randomized.PIT.isSubW, Randomized.PIT.isVarW}
For a natural number $k$, a function $a : br_{	ext{nat}} 	o br_{	ext{nat}}$, a list
$vs$ of natural numbers, and a Boolean list $w$, define its value by cases: for a
variable word, take the value of $a$ at the length of its first argument modulo $k$; for
a one word, take $1$ modulo $k$; for an addition word, add the entries of $vs$ indexed
by the lengths of its two arguments, using $0$ when an index is out of range, and reduce
modulo $k$; for a subtraction word, add the first such entry to $k$ minus the second and
reduce modulo $k$; in every other case, multiply the two entries and reduce modulo $k$.
\end{definition}

\begin{definition}[One step of the W process]
\lean{Randomized.PIT.wStep}
\label{Randomized.PIT.wStep}
\leanok
\uses{Randomized.PIT.isDone, Randomized.PIT.okW, Randomized.PIT.restW, Randomized.PIT.valN}
For natural numbers $n,k$, a function $a\colon\mathbb{N}\to\mathbb{N}$, and a state
$(b,v)$ consisting of a Boolean list and a natural-number list, define the next state as
follows: if $b$ is done, keep $(b,v)$ unchanged; otherwise, if $b$ is acceptable at
length $|v|$, replace it by its remainder and append the value determined by $n,k,a,v,b$
to $v$; if neither condition holds, return the empty Boolean list together with $v$.
\end{definition}

\begin{theorem}[Agreement of the Two Parsing Steps]
\lean{Randomized.PIT.absStep_eq_wStep}
\label{Randomized.PIT.absStep_eq_wStep}
\leanok
\uses{Complexity.BinArith.decodeNat_eq, Complexity.BinArith.hasFalse, Complexity.BinArith.uHead, Complexity.BinArith.uRest, Complexity.isNilB, Randomized.AlgGate.parse, Randomized.AlgGate.parseBin, Randomized.AlgGate.refsBelow, Randomized.PIT.absStep, Randomized.PIT.arg1W, Randomized.PIT.arg2W, Randomized.PIT.bodyW, Randomized.PIT.gateVal, Randomized.PIT.isAddW, Randomized.PIT.isDone, Randomized.PIT.isMulW, Randomized.PIT.isOneW, Randomized.PIT.isSubW, Randomized.PIT.isVarW, Randomized.PIT.midW, Randomized.PIT.okW, Randomized.PIT.rest2W, Randomized.PIT.restW, Randomized.PIT.tagBit, Randomized.PIT.valN, Randomized.PIT.wStep}
\difficulty{5}
Let $n,k$ be natural numbers, let $a\colon\mathbb{N}\to\mathbb{N}$, let $w$ be a finite
bit string, and let $v$ be a finite list of natural numbers. The transition that parses
an algebraic gate from a leading $1$ in $w$ and appends its value when all referenced
gates precede the current list length agrees with the transition that checks whether the
word is finished or valid and, when valid, appends its computed value: both produce the
same pair of residual word and value list from $(w,v)$.
\end{theorem}

\begin{theorem}[Equality of Two Circuit Evaluation Steps]
\lean{Randomized.PIT.wStep_eq_absStep}
\label{Randomized.PIT.wStep_eq_absStep}
\leanok
\uses{Randomized.PIT.absStep, Randomized.PIT.absStep_eq_wStep, Randomized.PIT.wStep}
\difficulty{4}
For every pair of natural numbers $n,k$ and every function
$a\colon\mathbb{N}\to\mathbb{N}$, the word-processing step on a pair consisting of a bit
string and a list of values agrees, on every such pair, with the abstract
gate-processing step using the same parameters.
\end{theorem}

\begin{definition}[First component of the context encoding]
\lean{Randomized.PIT.sCtx}
\label{Randomized.PIT.sCtx}
\leanok
\uses{Complexity.pairFstD}
For a Boolean list $D$, define its context component to be the first component extracted
from $D$.
\end{definition}

\begin{definition}[First component after two second-component extractions]
\lean{Randomized.PIT.sW}
\label{Randomized.PIT.sW}
\leanok
\uses{Complexity.pairFstD, Complexity.pairSndD}
For a Boolean list $D$, define the result by extracting the second component of $D$,
then the first component of that result, after first extracting the second component of
$D$ once. Equivalently, apply second-component extraction to $D$, then second-component
extraction again, and finally first-component extraction.
\end{definition}

\begin{definition}[First component after three second-component extractions]
\lean{Randomized.PIT.sI}
\label{Randomized.PIT.sI}
\leanok
\uses{Complexity.pairFstD, Complexity.pairSndD}
For a Boolean list $D$, define the result by applying second-component extraction three
times, then first-component extraction.
\end{definition}

\begin{definition}[Third successive second component]
\lean{Randomized.PIT.sL}
\label{Randomized.PIT.sL}
\leanok
\uses{Complexity.pairSndD}
For a Boolean list $D$, define the result by applying second-component extraction three
times.
\end{definition}

\begin{definition}[First component of the context component]
\lean{Randomized.PIT.sN}
\label{Randomized.PIT.sN}
\leanok
\uses{Complexity.pairFstD, Randomized.PIT.sCtx}
For a Boolean list $D$, define the result by taking its context component and extracting
the first component of that component.
\end{definition}

\begin{definition}[First component after context and second extraction]
\lean{Randomized.PIT.sRaw}
\label{Randomized.PIT.sRaw}
\leanok
\uses{Complexity.pairFstD, Complexity.pairSndD, Randomized.PIT.sCtx}
For a Boolean list $D$, define the result by taking its context component, extracting
the second component, and then extracting the first component.
\end{definition}

\begin{definition}[Second component after context and extraction]
\lean{Randomized.PIT.sPts}
\label{Randomized.PIT.sPts}
\leanok
\uses{Complexity.pairSndD, Randomized.PIT.sCtx}
For a Boolean list $D$, define the result by taking its context component, then
extracting the second component twice.
\end{definition}

\begin{definition}[Successor encoding of a Boolean list]
\lean{Randomized.PIT.sK}
\label{Randomized.PIT.sK}
\leanok
\uses{Complexity.BinArith.add, Randomized.PIT.sRaw, Turing.pairEncode}
For a Boolean list $D$, let $sK(D)$ be the binary addition of the pair encoding of the
raw representation of $D$ and the one-bit list $[1]$.
\end{definition}

\begin{definition}[Encoded word from two Boolean lists]
\lean{Randomized.PIT.xW}
\label{Randomized.PIT.xW}
\leanok
\uses{Complexity.BinArith.getW, Randomized.PIT.sL, Turing.pairEncode}
For Boolean lists $J$ and $D$, let $xW(J,D)$ be the word obtained by applying the
word-extraction operation to the pair encoding of $J$ and the list representation of
$D$.
\end{definition}

\begin{definition}[Raw value encoding]
\lean{Randomized.PIT.rawValW}
\label{Randomized.PIT.rawValW}
\leanok
\uses{Complexity.BinArith.add, Complexity.BinArith.getW, Complexity.BinArith.mul, Complexity.BinArith.sub, Randomized.PIT.arg1W, Randomized.PIT.arg2W, Randomized.PIT.isAddW, Randomized.PIT.isOneW, Randomized.PIT.isSubW, Randomized.PIT.isVarW, Randomized.PIT.sK, Randomized.PIT.sPts, Randomized.PIT.sW, Randomized.PIT.xW, Turing.pairEncode}
For a Boolean list $D$, define its raw value encoding by cases on its encoded word. In
the variable case, take the binary sum of the pair encoding of the extracted word at the
encoded variable index and the encoded point list; in the constant-one case, use the
one-bit list $[1]$. In the addition case, take the binary sum of the pair encoding of
the two recursively extracted operand words. In the subtraction case, take the binary
sum of the first operand word paired with the binary difference between the encoding of
$sK(D)$ paired with the second operand word. In every other case, take the binary
product of the pair encoding of the two extracted operand words.
\end{definition}

\begin{definition}[Reduced value encoding]
\lean{Randomized.PIT.valW}
\label{Randomized.PIT.valW}
\leanok
\uses{Complexity.BinArith.mod, Randomized.PIT.rawValW, Randomized.PIT.sK, Turing.pairEncode}
For a Boolean list $D$, define its value encoding as the binary remainder obtained from
the pair encoding of its raw value encoding and the encoding $sK(D)$.
\end{definition}

\begin{definition}[State transition encoding]
\lean{Randomized.PIT.gStep}
\label{Randomized.PIT.gStep}
\leanok
\uses{Randomized.PIT.isDone, Randomized.PIT.okW, Randomized.PIT.restW, Randomized.PIT.sCtx, Randomized.PIT.sI, Randomized.PIT.sL, Randomized.PIT.sN, Randomized.PIT.sW, Randomized.PIT.valW, Turing.pairEncode}
For a Boolean list $D$, define its next-step encoding as follows. If its encoded word is
done, retain $D$. Otherwise, if the validity condition holds for the lengths of its two
encoded lists and its encoded word, form the pair encoding of its context with the pair
encoding of the remaining word, the updated index list, and the concatenation of its
list representation with the encoding of its value. If that validity condition fails,
pair its context with the encoding of an empty word, the unchanged index list, and the
unchanged list representation.
\end{definition}

\begin{definition}[Context encoding]
\lean{Randomized.PIT.ctxOf}
\label{Randomized.PIT.ctxOf}
\leanok
\uses{Complexity.BinArith.encWords, Turing.pairEncode}
For a natural number $n$, a Boolean list $R$, and a list $P$ of Boolean lists, define
the context encoding as the pair encoding of the length-$n$ list of true values with the
pair encoding of $R$ and the word encoding of $P$.
\end{definition}

\begin{definition}[Point value]
\lean{Randomized.PIT.ptVal}
\label{Randomized.PIT.ptVal}
\leanok
\uses{Complexity.LogProg.bitsVal}
For a list $P$ of Boolean lists and an index $i$, define the point value to be one plus
the binary value of the $i$th list of $P$, using the empty list when that index is out
of range.
\end{definition}

\begin{definition}[State encoding]
\lean{Randomized.PIT.encState}
\label{Randomized.PIT.encState}
\leanok
\uses{Complexity.BinArith.encWords, Turing.pairEncode}
For Boolean lists $c$ and $w$ and a list $v$ of natural numbers, define their encoded
state by pairing $c$ with the result of pairing $w$ with the pair encoding of a list of
$|v|$ true values and the encoding of the list of binary representations of the entries
of $v$.
\end{definition}

\begin{theorem}[Encoded evaluator step]
\lean{Randomized.PIT.gStep_encState}
\label{Randomized.PIT.gStep_encState}
\leanok
\uses{Complexity.BinArith.add, Complexity.BinArith.encWords_append_singleton, Complexity.BinArith.getW_pairEncode, Complexity.BinArith.getW_pairEncode_bits, Complexity.BinArith.mod, Complexity.BinArith.mul, Complexity.BinArith.sub, Complexity.LogProg.bitsVal, Complexity.LogProg.bitsVal_bits, Complexity.pairFstD_pairEncode, Complexity.pairSndD_pairEncode, Randomized.PIT.ctxOf, Randomized.PIT.encState, Randomized.PIT.gStep, Randomized.PIT.ptVal, Randomized.PIT.rawValW, Randomized.PIT.sCtx, Randomized.PIT.sI, Randomized.PIT.sK, Randomized.PIT.sL, Randomized.PIT.sN, Randomized.PIT.sPts, Randomized.PIT.sRaw, Randomized.PIT.sW, Randomized.PIT.valN, Randomized.PIT.valW, Randomized.PIT.wStep, Randomized.PIT.xW}
\difficulty{5}
Let $n$ be a natural number, let $R$ and $w$ be finite bit strings, let $P$ be a finite
list of bit strings, and let $vs$ be a finite list of natural numbers. Encode the
context consisting of $n$, $R$, and $P$, together with $w$ and $vs$. One encoded
evaluator step on this state equals the encoding of the state obtained by one abstract
evaluator step on $(w,vs)$, using modulus equal to the binary value of $R$ plus one and
assigning input index $i$ the binary value of the $i$th entry of $P$ plus one (with a
missing entry treated as the empty bit string).
\end{theorem}

\begin{theorem}[Iteration of the gate-evaluation machine]
\lean{Randomized.PIT.iterate_gStep}
\label{Randomized.PIT.iterate_gStep}
\leanok
\uses{Complexity.LogProg.bitsVal, Randomized.PIT.absStep, Randomized.PIT.ctxOf, Randomized.PIT.encState, Randomized.PIT.gStep, Randomized.PIT.gStep_encState, Randomized.PIT.ptVal, Randomized.PIT.wStep_eq_absStep}
\difficulty{4}
Fix an input width $n$, a bit string $R$, and a list $P$ of bit strings. For every
number of steps $T$, work string $w$, and list of natural-number values $vs$, iterating
the encoded gate-evaluation machine for $T$ steps from the state with context determined
by $n,R,P$ and data $(w,vs)$ gives the encoded state with that same context and data
equal to the two components obtained by iterating the abstract decoding-and-evaluation
step for $T$ steps from $(w,vs)$, using modulus $\operatorname{val}(R)+1$ and point
assignment $i\mapsto \operatorname{val}(P_i)+1$.
\end{theorem}
